
Across the experiments, three empirical patterns were consistent. First, the baseline casewise derivative usually preserved the direction of change in the single-observation and batch settings. Second, the profiled-displacement surrogate \(Q_i^{\mathrm{disp}}\) was an inconsistent screen for finite-step failure. Third, stronger downweighting and multi-case paths often produced larger finite-step amplitudes than the first-order approximation predicted. Section~\ref{sec:higher_order_theory} gives a scalar profiled-parameter Poisson example in which \eqref{eq:ho_theta1} gives the local slope, \eqref{eq:ho_theta2} gives the second-order correction, and \eqref{eq:ho_theta_remainder} writes the finite-step miss as accumulated curvature.

The higher-order path figures clarify this pattern. In the generic directional analysis, \eqref{eq:ho_u1} and \eqref{eq:ho_u2} give the exact groupwise slope and curvature, and \eqref{eq:ho_u2_zero} shows that the inner second derivative vanishes only when the groupwise weighted residual aggregate is zero. The one-case bound \eqref{eq:ho_u2_lb} identifies the pointwise cases that bend most strongly: those with large \(|r_i|\) in weak-curvature groups. Equation \eqref{eq:ho_theta_remainder} gives the exact endpoint remainder, and \eqref{eq:ho_theta_remainder_lb} gives the corresponding local finite-step lower-bound mechanism for sufficiently small \(\tau\). This helps explain why the most separated pointwise cases were also the most strongly curved refit paths.

The profiled-displacement surrogate \(Q_i^{\mathrm{disp}}\) was less consistent as a screen for finite-step prediction error. This does not invalidate the original perturbation-selection literature, because the original score was designed for a different local target and a different normalization. In the present profiled-covariance problem, \(Q_i^{\mathrm{disp}}\) is a baseline Jacobian/Hessian surrogate, whereas finite-step behavior also depends on the second-order terms in \eqref{eq:ho_theta2}. Equation \eqref{eq:ho_theta2_zero} makes the consequence explicit: a flat path requires exact cancellation among those second-order terms.

The AMIP and random-subset controls show the same pattern for multi-case perturbations. Ranking cases by the local score usually produced a refit path larger than the predicted path. When the deleted sets were randomized, the gap remained in some datasets and changed sign in others. This means the finite response depended on how the perturbation set was assembled. In the Poisson example, \eqref{eq:ho_u2_zero} gives same-direction curvature for nonnegative downweighting directions whenever the relevant groupwise residual aggregate is nonzero, and the one-case corollary \eqref{eq:ho_theta_remainder_lb} shows why finite-step error is locally at least quadratic for sufficiently small \(\tau\) once that curvature survives the outer cancellation condition \eqref{eq:ho_theta2_zero}. For the outer covariance path, however, \eqref{eq:ho_theta2} still allows sign changes through \(F_{\vartheta\tau}\) and \(F_{\tau\tau}\). Thus the example gives a generic profiled-path mechanism, not a covariance-specific theorem.

The low-\(n/q\) pattern in the controlled Poisson random-intercept simulation matches both the mechanism in Section~\ref{sec:higher_order_theory} and the MCMC experiment in \citet[Sec.~4.1]{GiordanoBroderick2024BayesIJVariance}. In their Poisson random-effect experiment, the IJ covariance underestimates the simulated truth when \(N/G=1\) and improves markedly by \(N/G=10\), which they explain by poor concentration of the local random-effect posterior. Section~\ref{sec:higher_order_theory} gives the Laplace/MMLE analogue: the final low-occupancy corollary specializes the generic directional theory to the one-case path, \(H_i^{-1}\) is the local Laplace variance, and \eqref{eq:ho_curvature_ratio} shows that the curvature-to-slope ratio is order \(1/K = G/N\). Thus lower occupancy produces the same qualitative failure mode, weak local concentration and larger finite-step nonlinearity, even though our estimand is finite weighted-refit displacement of a profiled MMLE rather than the frequentist covariance of a posterior mean. Equation \eqref{eq:ho_theta_remainder_lb} also explains why, within one dataset, some observations separate much more than others at the same occupancy.

Four limitations matter. First, Section~\ref{sec:higher_order_theory} analyzes a scalar parameter of interest with many random effects. It captures the low-occupancy mechanism, but it is not intended to explain sign changes created by rotation of a genuinely multivariate covariance path under projection. Second, all derivatives are reported on the internal \texttt{lme4} covariance scale, so the paper studies local movement of that parameterization rather than a transformed substantive summary. Third, the derivative calculation assumes a fixed limiter state and is implemented only for Gaussian, binomial, and Poisson families, with REML sensitivity restricted to Gaussian models. If a perturbation crosses limiter boundaries or leaves that model class, explicit refits remain the relevant check. Fourth, the AMIP panels use \(\phi(\theta)=\|\theta\|_2\) only as a common scalar summary on the internal scale. They support statements about calibration and path shape, not strong comparisons of effect magnitude across datasets.

In practice, the derivative is best viewed as a fast screening tool rather than a substitute for explicit refitting. It provides a local sensitivity map, ranks perturbation directions, and supports structured path constructions such as the AMIP ranking. Explicit refits remain necessary for larger perturbations. Once path curvature, sign reversal, or limiter-state changes matter, the finite response can differ materially from the first-order prediction.
